%% kapa: exam=1 round=20 year=2009 period=2 form=A subject=accounting subject_ko=회계학 seq=17 number=17
%% source: Q-Net 감정평가사 기출문제 (cst00309) artlSeq=5024178 file=extracted/1008/5024178/2009년 제20회 감정평가사 1차시험 기출문제/2교시.hwp
%% conversion: kordoc 4.24.51 -> tools/gen1.py (auto); review=proofread-sonnet; flags=none
\begin{question}{17}
기초에서 기말까지 물가가 지속적으로 상승하고 있는 가운데 계속하여 매출이 이루어지고 기초에 비하여 기말재고자산의 수량이 감소하지 않은 경우, 매출총이익을 가장 크게 표시하는 재고자산의 원가결정방법부터 순서대로 나열한 것은? (단, 재고청산이 없다고 가정한다.)
\begin{choices}
  \item 선입선출법-이동평균법-총평균법-후입선출법
  \item 선입선출법-총평균법-이동평균법-후입선출법
  \item 총평균법-선입선출법-이동평균법-후입선출법
  \item 이동평균법-선입선출법-총평균법-후입선출법
  \item 선입선출법-후입선출법-총평균법-이동평균법
\end{choices}
\end{question}
\answer{1}
